\section{Proofs for Gaussian binary least squares}
\label{app:bls-proofs}

\subsection{Elementary concentration estimates}
\label{app:bls-probability-lemmas}

We use the following estimates at finite dimensions.  If $G$ is standard
normal, $Q_d\sim\chi^2_d$, and $K_n\sim\operatorname{Bin}(n,1/2)$, then
for $z,x,s\ge0$,
\begin{align}
 \Pbb\{|G|>z\}&\le e^{-z^2/2},
 \label{app:bls-normal-tail}\\
 \Pbb\{Q_d>d+2\sqrt{dx}+2x\}&\le e^{-x},\qquad
 \Pbb\{Q_d<d-2\sqrt{dx}\}\le e^{-x},
 \label{app:bls-chi-square-tails}\\
 \Pbb\{|K_n-n/2|>s\}&\le2e^{-2s^2/n}.
 \label{app:bls-binomial-tail}
\end{align}
The convention $Q_0=0$ makes the second line valid for zero degrees.
For completeness, the first inequality follows from
$2e^{z^2/2}\overline\Phi(z)\le1$: the expression equals one at zero
and is decreasing, since
$z\overline\Phi(z)\le(2\pi)^{-1/2}e^{-z^2/2}$ by integrating
$t e^{-t^2/2}$ over $t\ge z$.
The Gaussian integral gives
$\E e^{tQ_d}=(1-2t)^{-d/2}$ for $t<1/2$.
For $0<t<1/2$ the upper centered moment satisfies
\[
 \log\E e^{t(Q_d-d)}\le\frac{dt^2}{1-2t}.
\]
For every $t\ge0$ the lower centered moment satisfies
\[
 \log\E e^{-t(Q_d-d)}\le dt^2.
\]
The inequalities follow by expanding $-\log(1-2t)$ and by
$\log(1+2t)\ge2t-2t^2$, respectively.  Exponential Markov bounds,
with $t=\sqrt{x/d}/(1+2\sqrt{x/d})$ for the upper tail and
$t=\sqrt{x/d}$ for the lower tail, give
\eqref{app:bls-chi-square-tails} for $d>0$; the case $d=0$ is immediate.
Finally,
$\E e^{t(K_n-n/2)}=\cosh(t/2)^n\le e^{nt^2/8}$; optimizing Markov's
inequality gives \eqref{app:bls-binomial-tail}.

\begin{lemma}[Concentration of the projection mixture]
\label{app:bls-beta-concentration}
Let $X\sim\chi^2_s$ and $Z\sim\chi^2_{M-s}$ be independent and
$Y=X/(X+Z)$, where $0\le s\le M$ and $M\ge1$.
For $0\le t\le1$,
\[
 \Pbb\{|Y-s/M|>t\}\le4e^{-Mt^2/64}.
\]
If instead $s\sim\operatorname{Bin}(n,1/2)$, $n\le M$, and $X,Z$
have these laws conditional on $s$, then
\[
 \Pbb\{|Y-n/(2M)|>2t\}\le6e^{-Mt^2/64}.
\]
\end{lemma}

\begin{proof}
The statement is trivial at $t=0$.  Put $x=Mt^2/64$.  Except on an event
of probability at most $2e^{-x}$, the chi-square estimates give
\[
 X\le s+\delta_1,\qquad Z\ge M-s-\delta_2,
 \quad \delta_1=2\sqrt{sx}+2x,
 \quad \delta_2=2\sqrt{(M-s)x}.
\]
Here $\delta_1\le M(t/4+t^2/32)$ and $\delta_2\le Mt/4$.
Monotonicity of the ratio, or the bound $Y\le1$ when the displayed lower
bound for $Z$ is negative, gives
\[
 Y-\frac sM
 \le\frac{(M-s)\delta_1+s\delta_2}
          {M(M+\delta_1-\delta_2)}
 \le\frac{t/4+t^2/32}{1-t/4}\le\frac{3t}{8}<t.
\]
Interchanging $X,Z$ proves the lower tail with the same failure bound.
For the mixture, \eqref{app:bls-binomial-tail} bounds
$\Pbb\{|s-n/2|>Mt\}$ by
$2e^{-2M^2t^2/n}\le2e^{-2Mt^2}$.  On the complement,
$|s/M-n/(2M)|\le t$.  Combining the two estimates proves the result.
\end{proof}

\begin{lemma}[An elementary Gaussian operator bound]
\label{app:bls-operator-bound}
Suppose $M_N/N\to\beta<\infty$, and define
\begin{equation}
 C_\beta=2\left(\sqrt{\beta+1}+
                  \sqrt{2(\log5+1)}\right).
 \label{app:bls-operator-constant}
\end{equation}
For an $M_N\times N$ independent standard Gaussian matrix $H_N$,
$\Pbb\{\|H_N\|_{\mathrm{op}}>C_\beta\sqrt N\}\le e^{-N}$ for all
sufficiently large $N$.
\end{lemma}

\begin{proof}
A maximal $1/2$-separated subset $\mathcal V$ of the unit sphere in
$\R^N$ is a $1/2$-net of size at most $5^N$.  Indeed, its disjoint
radius-$1/4$ balls lie in the radius-$5/4$ ball; comparing volumes gives
the bound.  For any matrix,
$\|H\|_{\mathrm{op}}\le2\max_{z\in\mathcal V}\|Hz\|$ by approximating
a maximizing unit vector by a net point.  Each $\|Hz\|^2$ has law
$\chi^2_M$.  Apply \eqref{app:bls-chi-square-tails} with
$x=(\log5+1)N$ and note that
$M+2\sqrt{Mx}+2x\le(\sqrt M+\sqrt{2x})^2$.
The union bound leaves probability at most $e^{-N}$.  Eventually
$M\le(\beta+1)N$, giving the stated constant.
\end{proof}

\subsection{Recovery and the value gap}
\label{app:bls-ml-proof}

We prove \Cref{bls:ml}.  For a deterministic $c\in\R^N$,
$Bc\sim N(0,a_c^2I_M)$ independently of $w$, where
$a_c^2=\rho\|c\|^2/N$.  Direct integration over one coordinate gives
\[
 \E e^{-\lambda(F(c)-W)}
 =\left[1+(2\lambda-4\lambda^2)a_c^2\right]^{-M/2}
 \quad(0<\lambda<1/2).
\]
For example, conditional on $Bc=z$, integration over $w$ gives
$\exp((2\lambda^2-\lambda)\|z\|^2)$, and the remaining Gaussian
integral gives the formula.  At $\lambda=1/4$, exponential Markov gives
\begin{equation}
 \Pbb\{F(c)\le W-s\}
 \le e^{-s/4}\left(1+\frac{\rho\|c\|^2}{4N}\right)^{-M/2}
 \quad(s\ge0).
 \label{app:bls-direction-chernoff}
\end{equation}
At a vertex $c=2\mathbf1_S$.  Hence
\begin{equation}
 \Pbb\{\exists S\ne\varnothing:F(2\mathbf1_S)\le W-s\}
 \le e^{-s/4}\Sigma_N,\qquad
 \Sigma_N=\sum_{k=1}^N\binom Nk(1+\rho k/N)^{-M/2}.
 \label{app:bls-vertex-sum}
\end{equation}

First consider recovery.  Since every term of $\Sigma_N$ decreases in
$\rho$, it suffices to set
$\rho=(2+\delta)\log N/\beta_N$, reducing $\delta$ if necessary so
$0<\delta\le1$.  Write $x_k=\rho k/N$ and set $x_0=\delta/4$.
If $x_k\le x_0$, then
$\log(1+x_k)\ge x_k(1-x_0/2)$, and
\[
 \binom Nk(1+x_k)^{-M/2}
 \le\frac{[N e^{-(\beta_N\rho/2)(1-\delta/8)}]^k}{k!}
 \le\frac{N^{-\delta k/4}}{k!}.
\]
The sum of these terms tends to zero.  For $x_k>x_0$, put $t=k/N$.
Using $\binom Nk\le(eN/k)^k$, the logarithm of the term divided by $N$
is at most
\[
 t\log(e/t)-\frac{\beta_N}2\log(1+\rho t).
\]
Choose $t_1>0$ so small that
$t_1\log(e/t_1)\le\tfrac14\log(1+x_0)$.
For $t\le t_1$ the display is at most
$-\tfrac14\log(1+x_0)$, since $\beta_N\ge1$.
For $t\ge t_1$ it is at most
$1-\tfrac12\log(1+\rho t_1)$, which tends to $-\infty$.
The large-$x_k$ sum is therefore $O(Ne^{-cN})$ for some $c>0$.
With $s=0$, \eqref{app:bls-vertex-sum} shows that no other vertex has
value at most $W$ with probability tending to one.  This proves uniqueness.

For the uniform value comparison, split at $x_k=7$ instead.  On $[0,7]$,
concavity gives $\log(1+x)\ge(\log8)x/7>0.28x$.  Thus the small-$x_k$
sum is at most
\[
 \sum_{k=1}^N\binom Nk e^{-0.14\beta_N\rho k}
 =(1+e^{-0.14\beta_N\rho})^N-1.
\]
For $x_k>7$, the same entropy bound gives exponent at most
$N(1-\tfrac12\log8)=-c_0N$, where $c_0>0$.  Set
$s=4N\log(1+e^{-0.14\beta_N\rho})+4\log N$ in
\eqref{app:bls-vertex-sum}.  The probability is at most
$N^{-1}+Ne^{-c_0N}$, proving \eqref{bls:ml-value-gap}.
The binomial expression, rather than its looser exponential upper bound,
is needed for the displayed logarithmic gap.

\subsection{The two root exactness probabilities}
\label{app:bls-root-exactness-proof}

We prove \Cref{bls:root-exactness}.  Put $\zeta_i=b_i^Tw$.
The first-order conditions for $u=0$ to minimize the convex function
$F$ on $[0,2]^N$ are $\zeta_i\ge0$ for all $i$.
Conditional on $w\ne0$, these inner products are independent centered
nondegenerate Gaussian variables.  Thus their signs are independent fair
coins, proving \eqref{bls:planted-root-probability}.  The matrix has full
column rank almost surely, making the objective strictly convex.  A root
value equal to the integer optimum therefore forces the unique root
minimizer to be an optimal vertex, and a vertex root minimizer conversely
forces value equality.

Fix a flipped vertex indexed by $S$, $|S|=k$, and write
$z=B_S\mathbf1$ and $v=w+2z$.  Its box optimality conditions are
$b_i^Tv\ge0$ off $S$ and $-b_i^Tv\ge0$ on $S$.
Given $(w,B_S)$, the $N-k$ off-support conditions hold with probability
$2^{-(N-k)}$.  Summing the conditions on $S$ gives the necessary condition
$-z^Tw-2\|z\|^2\ge0$.  Here $z\sim N(0,a^2I_M)$ independently of
$w$, with $a^2=\rho k/N$.  For $\lambda>0$ in the range where the moment
exists,
\[
 \E e^{\lambda(-z^Tw-2\|z\|^2)}
 =\left(1+4\lambda a^2-\lambda^2a^2\right)^{-M/2}.
\]
At $\lambda=2$, Markov's inequality bounds the necessary event by
$(1+4\rho k/N)^{-M/2}$.  Its indicator is measurable with respect to
$(w,B_S)$, so multiplication by the conditional off-support probability
is valid.  Summing over vertices gives the first inequality in
\eqref{bls:global-root-probability}; for $k=0$ use the already proved
probability $2^{-N}$.
Finally,
$\log(1+4\rho k/N)\ge(k/N)\log(1+4\rho)$ by concavity.  The binomial
theorem then gives the second inequality.

\subsection{Coefficient density and projection mixture}
\label{app:bls-nnls-proofs}

We prove \Cref{bls:nnls-law}.  Full column rank and the nonzero
conditions of \Cref{bls:sign-orbit} hold almost surely: rank-deficient
Gaussian matrices have Lebesgue measure zero, and, conditional on a
full-rank matrix, each required functional of $v$ is nondegenerate
Gaussian.  The sign-orbit identity gives a uniform support and hence
$K\sim\operatorname{Bin}(n,1/2)$.

Fix a nonempty support $S$ of size $s$.  Conditional on $H_S$,
\[
 d^S\sim N\left(0,\frac{c^2}{a^2}(H_S^TH_S)^{-1}\right).
\]
Multiplying its conditional density by the Gaussian matrix density shows
that its unconditional density at $d$ is proportional to
\[
 \left(\frac ac\right)^s
 \int_{\R^{M\times s}}\det(H^TH)^{1/2}
 \exp\left[-\frac12\operatorname{tr}
       \left(H^TH\left(I_s+\frac{a^2}{c^2}dd^T\right)\right)\right]dH.
\]
Put $T=I_s+(a^2/c^2)dd^T$ and change variables $H=GT^{-1/2}$.
The transformation of the $M$ rows has Jacobian $\det(T)^{-M/2}$;
the Gram determinant contributes $\det(T)^{-1/2}$.  The exponential
becomes $\exp(-\operatorname{tr}(G^TG)/2)$.  Thus all dependence on $d$
is
\[
 \det(T)^{-(M+1)/2}
 =\left(1+\frac{a^2}{c^2}\|d\|^2\right)^{-(M+1)/2}.
\]
The matrix integral is finite, since Hadamard's inequality bounds
$\det(G^TG)^{1/2}$ by the product of the column norms, whose Gaussian
moments are finite.  This argument also covers $M=s$.

To compute the normalizing constant, take independent
$Z_s\sim N(0,I_s)$ and $Q\sim\chi^2_\nu$, with $\nu=M-s+1$.
For $D=(c/a)Z_s/\sqrt Q$, condition on $Q=q$ and change variables in the
Gaussian density.  Integration over $q$ gives
\begin{align*}
 p_D(d)
 &=\frac{(a/c)^s}{2^{(s+\nu)/2}\pi^{s/2}\Gamma(\nu/2)}
   \int_0^\infty q^{(s+\nu)/2-1}
   e^{-q(1+(a^2/c^2)\|d\|^2)/2}\,dq\\
 &=\left(\frac ac\right)^s
 \frac{\Gamma((s+\nu)/2)}{\pi^{s/2}\Gamma(\nu/2)}
 \left(1+\frac{a^2}{c^2}\|d\|^2\right)^{-(s+\nu)/2}.
\end{align*}
Since $s+\nu=M+1$, this is \eqref{bls:student-density}.  It is a
normalized density with the shape just derived for $d^S$, so the two laws
agree.

For any Borel set $D\subseteq[0,\infty)^s$, the event
$|d^S|\in D$ is invariant under every column sign change.
\Cref{bls:sign-orbit}, divided by $\Pbb\{\mathcal S=S\}=2^{-n}$,
therefore gives
\[
 \mathcal L\bigl((u_j^\circ)_{j\in S}\mid\mathcal S=S\bigr)
 =\mathcal L(|d^S|).
\]
This proves \eqref{bls:coefficient-mixture}.  Conditional on $Q$,
the probability that its $s$ absolute normal coordinates are all below
$at\sqrt Q/c$ is
$[1-2\overline\Phi(at\sqrt Q/c)]^s$.
Mixing over the binomial support size gives \eqref{bls:maximum-cdf}.
At $t=0$ every positive-size summand vanishes, leaving the atom $2^{-n}$;
the explicitly defined zero-size summand avoids an ambiguous $0^0$.

For the projection statement, fix $S$ and set
$X_S=\|P_{L_S}v\|^2$ and $Z_S=\|(I-P_{L_S})v\|^2$.
These quantities are invariant under all column sign changes.  Conditional
on $C$, their scaled laws are independent $\chi^2_s$ and
$\chi^2_{M-s}$; these laws depend only on the dimensions.
The invariant-event identity transfers their joint distribution to the
event $\mathcal S=S$.  On that event the fitted vector is
$Cu^\circ=-P_{L_S}v$ and the residual is $(I-P_{L_S})v$.
Thus the same joint law holds for $X,Z$ conditional on the selected
support.  Supports of the same size have identical laws, proving the
claimed mixture and \eqref{bls:projection-concentration} by
\Cref{app:bls-beta-concentration}.

This proof transfers sign-invariant quantities from a fixed support.
It does not treat the signed regression vector as independent of the
event selecting its support.

\subsection{The matching root inactivity scale}
\label{app:bls-root-threshold-proof}

Here we complete the asymptotic argument in \Cref{bls:root-threshold}.
Use its synthetic representation with
$K\sim\operatorname{Bin}(N,1/2)$,
$Q\mid K\sim\chi^2_{M-K+1}$, and conditionally independent normal
coordinates $Z_j$.
\eqref{app:bls-binomial-tail} gives $K/N\to1/2$ in probability and
$\Pbb\{N/4\le K\le3N/4\}\to1$.
On this event the chi-square degree $d=M-K+1$ is between $N/4$ and
$(\beta+1)N+1$ for all sufficiently large $N$.
For each fixed $\eta>0$, choosing a sufficiently small constant $c_\eta>0$
and $x=c_\eta N$ in \eqref{app:bls-chi-square-tails} bounds
$\Pbb\{|Q-d|>\eta N\mid K\}$ by $2e^{-c_\eta N}$ uniformly on
that event.  Therefore
\[
 \frac QN-\left(\beta_N-\frac12\right)
 =\frac{Q-(M-K+1)}N-\left(\frac KN-\frac12\right)+\frac1N
 \xrightarrow{\Pbb}0.
\]

For $\delta\in(0,1)$, the Gaussian tail bound and $K\le N$ give
\[
 \Pbb\{\max_{j\le K}|Z_j|^2>2(1+\delta)\log N\}
 \le N^{-\delta}.
\]
For the reverse inequality, integrate the normal density over
$[z,z+1/z]$ for $z\ge1$.  The density throughout this interval is at
least a fixed positive constant times $e^{-z^2/2}$, so
$\Pbb\{|G|>z\}\ge c z^{-1}e^{-z^2/2}$.
At $z=\sqrt{2(1-\delta)\log N}$, this yields
\[
 \Pbb\{|G|>z\}\ge
 c_\delta\frac{N^{-(1-\delta)}}{\sqrt{\log N}}.
\]
Conditional on $K\ge N/4$, independence of the normal coordinates then
bounds the probability of no exceedance by
\[
 \exp\left[-\frac{c_\delta N^\delta}{4\sqrt{\log N}}\right].
\]
It follows that $\max_{j\le K}|Z_j|^2/(2\log N)\to1$ in probability.
Since $\beta_N-1/2\ge1/2$, the denominator $Q/N$ stays bounded away
from zero in probability, and the ratio in
\eqref{bls:root-cutoff-scale} follows.

For completeness, when $\rho=O(\log N)$ the logarithm of
\eqref{bls:root-inactivity-tail} has the asserted precision.  Taylor
expansion gives
\[
 \log q=-2\rho/N+O(\rho^2/N^2),\qquad
 \log\frac{1+q}{2}=-\rho/N+O(\rho^2/N^2).
\]
Multiplying by $M-N+1$ and $N-1$, respectively, gives
\[
 \log\left[\frac N2q^{M-N+1}
                 \left(\frac{1+q}{2}\right)^{N-1}\right]
 =\log(N/2)-(2\beta_N-1)\rho
    +O(\rho/N+\rho^2/N).
\]
Boundedness of $\beta_N$ and $\rho=O(\log N)$ make the remainder $o(1)$.

\subsection{The static variable-order upper bound}
\label{app:bls-static-upper-proof}

We prove the upper bound in \Cref{bls:certificate-law} directly from the
orthant residual mixture.  Assume $\rho_N\to\infty$ and $\rho_N=o(N)$,
and abbreviate $\rho=\rho_N$.  A permutation fixed independently of
the data can be taken as the identity by Gaussian column exchangeability.
Define
\[
 a_N=\frac{\log(e\rho)}\rho+\frac{\log N}N+\frac\rho N,
 \qquad \eta_N=a_N^{1/8},\qquad
 K=\left\lceil\frac{(1+\eta_N)N}{4(2\beta_N-1)\rho}\right\rceil.
\]
Then $\eta_N\to0$, $K\to\infty$, $K/N\to0$, and
\begin{equation}
 \log\binom NK\le K\log(eN/K)=O(N\log(e\rho)/\rho)=o(N).
 \label{app:bls-small-set-entropy}
\end{equation}
Put $L_N=\log(2(N+1)\binom NK)+2\log N$, so $L_N/M=O(a_N)$.
At a depth-$d$ node with wrong set $T\subseteq[d]$ of size $K$, the
free set is $\{d+1,\ldots,N\}$.  Its residual is independent of its
free matrix and has law $N(0,c_K^2I_M)$, where
$c_K^2=1+4\rho K/N$.  This variance stays bounded because
$\rho K/N=(1+\eta_N)/[4(2\beta_N-1)]+O(\rho/N)$.
Writing $\widetilde r_{d,T}$ for the orthant value,
\Cref{bls:nnls-law} gives its marginal representation
\[
 \widetilde r_{d,T}/c_K^2\mid J_{d,T}=s\sim\chi^2_{M-s},
 \qquad J_{d,T}\sim\operatorname{Bin}(N-d,1/2).
\]
This also includes no free coordinates by taking $J_{d,T}=0$.

For a fixed node, the binomial upper tail gives
$J_{d,T}\le(N-d)/2+\sqrt{NL_N}$ except on probability at most
$e^{-2L_N}$.  Conditional on this index, the chi-square lower tail gives
$\widetilde r_{d,T}/c_K^2\ge M-J_{d,T}-2\sqrt{ML_N}$ except on
probability at most $e^{-L_N}$; replacing its degree by $M$ in the square
root weakens the lower bound.  There are at most
$(N+1)\binom NK$ such fixed-order nodes.  A union bound therefore proves
that, with probability tending to one, simultaneously
\[
 \widetilde r_{d,T}
 \ge c_K^2\left(M-\frac{N-d}2\right)-O(\sqrt{NL_N})
 \ge c_K^2(M-N/2)-O(\sqrt{NL_N}).
\]
The box value is at least the orthant value.  Moreover,
\[
 c_K^2(M-N/2)-M
 =2\rho K(2\beta_N-1)-\frac N2\ge\frac{\eta_NN}{2}.
\]
The error is $O(N\sqrt{a_N})$, while
$W=M+O_{\Pbb}(\sqrt{N\log N})$.  Since
$\sqrt{a_N}=o(\eta_N)$, every $K$-wrong node has box bound at least
$W$ with probability tending to one.  It is pruned for any incumbent of
value at most $W$.  Dependence between these node values and $W$ does
not matter: the events are intersected by a union bound.

No node with more than $K$ wrong fixings is generated.  Every node that
branches has at most $K-1$ wrong fixings and depth less than $N$.
Counting such potential parents, including ones that would already have
been pruned, gives
\[
 \#\text{nodes}\le
 1+2\sum_{d=0}^{N-1}\sum_{j=0}^{K-1}\binom dj
 =1+2\sum_{i=1}^K\binom Ni.
\]
For $K\le N/2$, put $t=K/N$.  Since $t^{-i}\le t^{-K}$ for $i\le K$,
\[
 \sum_{i=0}^K\binom Ni
 \le t^{-K}(1+t)^N\le(eN/K)^K.
\]
This proves \eqref{bls:static-tree-count}.  As $K\to\infty$,
rounding has relative error $o(1)$ and
\[
 K\log(eN/K)
 =(1+o(1))\frac{N}{4(2\beta-1)\rho}\log\rho.
\]
The logarithm of the node count has the same upper bound because the
last expression diverges.  With an optimal incumbent, the leaves certify
$\OPT$, and hence $\OPT-\varepsilon_N$.  This establishes the needed
upper certificate order without assuming $\OPT=W$.

\subsection{An entropy bound for overlapping supports}
\label{app:bls-entropy-proof}

\begin{lemma}[Overlap forces an entropy loss]
\label{app:bls-overlap-entropy}
Let $\mathcal F$ be a nonempty family of $k$-subsets of an $n$-set, where
$1\le k<n$, and let $p_i$ be the fraction of its members containing $i$.
Suppose $0<\lambda\le1$,
$\sum_i p_i^2\ge\lambda k$, and
$\lambda/2\ge e^2k/n$.  Then
\begin{equation}
 \log\frac{\binom nk}{|\mathcal F|}
 \ge\frac{\lambda k}{2}\log\frac n{ek}-\log(8\sqrt k).
 \label{app:bls-entropy-loss}
\end{equation}
\end{lemma}

\begin{proof}
Let $S$ be uniform on $\mathcal F$, represented by its indicator vector.
Entropy subadditivity gives
$\log|\mathcal F|\le\sum_i h(p_i)$, where
$h(p)=-p\log p-(1-p)\log(1-p)$, with $0\log0=0$.
Indeed, the entropy chain rule writes the joint entropy as a sum of
conditional entropies, and conditioning cannot increase entropy.
For $q=k/n$, $\sum_i p_i=k$ implies
\[
 \sum_i h(p_i)=nh(q)-\sum_i d(p_i\Vert q),
\]
where $d(p\Vert q)=p\log(p/q)+(1-p)\log((1-p)/(1-q))$.
Using $\log x\ge1-1/x$ on the second term gives
\[
 d(p\Vert q)\ge\phi(p):=p\log(p/q)+q-p\ge0.
\]
The last inequality follows since $\phi$ is convex with minimum zero at
$p=q$.  Put $L=\log(1/(eq))\ge1$.
We claim $\phi(p)\ge Lp(p-\lambda/2)$ for $0\le p\le1$.
For $p\le\lambda/2$ this follows from nonnegativity.
For $p\ge\lambda/2$ it suffices that
$\log(p/q)-1-L(p-\lambda/2)\ge0$.
This function is concave; at $p=\lambda/2$ it equals
$\log(\lambda/(2q))-1\ge1$, and at $p=1$ it equals
$L\lambda/2\ge0$.  The claim follows throughout the interval.
Summing gives
\[
 \sum_i d(p_i\Vert q)
 \ge L\left(\sum_i p_i^2-\frac\lambda2\sum_i p_i\right)
 \ge\frac{L\lambda k}{2}.
\]
Thus $\log|\mathcal F|\le nh(q)-L\lambda k/2$.

We also need a lower estimate for $\binom nk$.
For $V\sim\operatorname{Bin}(n,q)$, $k=nq$ is a mode, as follows by
taking ratios of consecutive probabilities.  Since
$\operatorname{Var}(V)=nq(1-q)\le k$, Chebyshev's inequality places
probability at least $1/2$ on $|V-k|<\sqrt{2k}$.
This interval contains at most $4\sqrt k$ integers, each with probability
at most $\Pbb\{V=k\}$.  Therefore
\[
 \binom nk q^k(1-q)^{n-k}=\Pbb\{V=k\}
 \ge\frac1{8\sqrt k}.
\]
Taking logarithms gives
$\log\binom nk\ge nh(q)-\log(8\sqrt k)$.
Combining the estimates proves \eqref{app:bls-entropy-loss}.
\end{proof}

\subsection{The convex-piece lower bound}
\label{app:bls-hard-lower-proof}

We prove \eqref{bls:finite-hard-lower} in precisely the two regimes used
in \Cref{bls:certificate-law}: divergent $\rho_N=o(N)$, or fixed
sufficiently large $\rho$.  Constants here are intentionally
conservative, allowing an elementary operator estimate.
Set
\begin{equation}
 p=\Phi(-1)-\Phi(-2)>0,\qquad
 d_\beta=\frac{\beta p}{64e^2C_\beta^2},\qquad
 c_\beta=\frac{\beta p}{256e^2C_\beta^4},
 \label{app:bls-hard-constants}
\end{equation}
where $C_\beta$ is defined in \eqref{app:bls-operator-constant}.
Take $\rho_0\ge e$ large enough that
\begin{equation}
 4\log(1+e^{-0.07\beta\rho_0})<d_\beta/2.
 \label{app:bls-rho-zero-choice}
\end{equation}
The fixed-SNR conclusion uses this choice and not a numerical threshold
from a sharper spectral estimate.

For $w\ne0$, decompose each signed standard Gaussian column as
\[
 \widetilde h_i=g_i\widehat w+h_i^\perp,
 \qquad \widehat w=w/\|w\|,
 \quad g_i=\widetilde h_i^T\widehat w.
\]
Conditional on $w$, the $g_i$ are independent standard normals, and their
orthogonal components are independent Gaussian vectors in
$\widehat w^\perp$, independent of the $g_i$.
Let
\[
 T=\{i:-2\le g_i\le-1\},\qquad n'=|T|,
 \qquad a=\sqrt{\rho/N}\,\|w\|.
\]
With probability tending to one, $n'\ge pN/2$, and on $T$,
$a\le-\zeta_i\le2a$, where $\zeta_i=b_i^Tw=ag_i$.
Binomial concentration proves the first assertion; it is valid after
conditioning on $w$ because the conditional law does not depend on $w$.
By the chi-square estimates, also with probability tending to one,
\begin{equation}
 \frac{\beta N}{2}\le W\le2\beta N,
 \qquad \frac12\sqrt{\beta\rho}\le a\le2\sqrt{\beta\rho}.
 \label{app:bls-noise-good-event}
\end{equation}
The coarse factors cover $\beta_N\to\beta$.
Let $H^\perp$ be the matrix of all orthogonal components.  The
deterministic inequality
$\|H^\perp\|_{\mathrm{op}}\le\|\widetilde H\|_{\mathrm{op}}$ and
\Cref{app:bls-operator-bound} give, on another event of probability
tending to one,
\begin{equation}
 \|H^\perp\|_{\mathrm{op}}\le C_\beta\sqrt N.
 \label{app:bls-perpendicular-operator}
\end{equation}
This uses the full matrix, so no conditioning on the selected set $T$ is
needed.  We intersect these events by a union bound.

Put
\begin{equation}
 \lambda=\frac{\sqrt\beta}{4C_\beta^2\sqrt\rho},\qquad
 k=\left\lfloor\frac{\lambda n'}{2e^2}\right\rfloor,
 \qquad
 P'=\{x^S:S\subseteq T,\ |S|=k\}.
 \label{app:bls-selected-vertices}
\end{equation}
In either stated regime, $k\to\infty$, so eventually $1\le k<n'$.
Also $\lambda\le1$ and $k/n'\le\lambda/(2e^2)$.
The lower bound on $n'$ and \eqref{app:bls-noise-good-event} imply
\begin{align*}
 ak&\ge\frac12\sqrt{\beta\rho}
       \left(\frac{\lambda pN}{4e^2}-1\right)\\
   &=\frac{\beta pN}{32e^2C_\beta^2}
       -\frac12\sqrt{\beta\rho}\ge d_\beta N
\end{align*}
for all sufficiently large $N$.  The final inequality uses
$\sqrt\rho=o(N)$, which holds in both regimes.

Let $\Delta=W-\OPT$.  In the divergent regime,
\Cref{bls:ml} gives $\Delta=o(N)$ with probability tending to one;
also $\varepsilon_N\le\rho_N=o(N)$.
At a fixed $\rho\ge\rho_0$, eventually $\beta_N\ge\beta/2$, and
\eqref{bls:ml-value-gap} and \eqref{app:bls-rho-zero-choice} give
$\Delta<d_\beta N/2+4\log N$ with probability tending to one, while
$\varepsilon_N\le\rho$ is bounded.  Consequently, in either regime,
\begin{equation}
 \Delta+\varepsilon_N\le ak
 \label{app:bls-gap-paid}
\end{equation}
holds with probability tending to one.  All the following statements are
deterministic on the intersection of these events.

Consider any class $I$ admissible for the target
$\OPT-\varepsilon_N$, and suppose $I\cap P'$ is nonempty.
Let $\mathcal F_I=\{S:x^S\in I\cap P'\}$ and let $p_i$ be the
fraction of these sets containing $i$; set $p_i=0$ outside $T$.
Then $\sum_i p_i=k$, and the barycenter of these vertices has error
coordinates $c_i=2p_i$.  It belongs to $\operatorname{conv} I$.
Define $X=-\zeta^Tc$.  Since all selected coordinates lie in $T$,
$2ak\le X\le4ak$.
The orthogonal decomposition gives the exact identity
\begin{equation}
 F(c)-W=-2X+\frac{X^2}{W}
                  +\frac\rho N\|H^\perp c\|^2.
 \label{app:bls-barycenter-identity}
\end{equation}
Indeed, the component of $Bc$ parallel to $w$ is
$(\zeta^Tc/\|w\|)\widehat w$.
The choices above yield
\[
 \frac XW\le\frac{4ak}{W}
 \le\frac{2}{e^2C_\beta^2}<\frac12.
\]
Thus $-2X+X^2/W\le-3X/2\le-3ak$.
Also \eqref{app:bls-perpendicular-operator} bounds the last term in
\eqref{app:bls-barycenter-identity} by
$4\rho C_\beta^2\sum_i p_i^2$.
Admissibility forces $F(c)\ge\OPT-\varepsilon_N
=W-\Delta-\varepsilon_N$, and therefore
\begin{equation}
 \sum_i p_i^2
 \ge\frac{3ak-\Delta-\varepsilon_N}{4\rho C_\beta^2}
 \ge\frac{ak}{2\rho C_\beta^2}
 \ge\lambda k.
 \label{app:bls-forced-overlap}
\end{equation}
The last step uses $a\ge\tfrac12\sqrt{\beta\rho}$.
Every admissible class hence meets $P'$ in a family satisfying
\Cref{app:bls-overlap-entropy}.

If $m$ is the largest possible size of such an intersection, then covering
$P'$ requires at least $\binom{n'}k/m$ classes.  Applying
\eqref{app:bls-entropy-loss} gives
\begin{equation}
 \log\kappa_{\OPT-\varepsilon_N}
 \ge\frac{\lambda k}{2}\log\frac{n'}{ek}
                         -\log(8\sqrt k).
 \label{app:bls-class-count}
\end{equation}
Now
\[
 \log\frac{n'}{ek}\ge\log\frac{2e}\lambda\ge\frac12\log\rho,
 \qquad
 \lambda k\ge\frac{\lambda^2pN}{4e^2}-\lambda
 =\frac{\beta pN}{64e^2C_\beta^4\rho}-\lambda.
\]
Substitution gives
\[
 \log\kappa_{\OPT-\varepsilon_N}
 \ge c_\beta\frac N\rho\log\rho
       -\frac\lambda4\log\rho-\log(8\sqrt k).
\]
The quantity $\lambda\log\rho$ is bounded for $\rho\ge\rho_0$.
Since $k\le N$, the two negative terms are at most $\log(8N)$ for all
sufficiently large $N$.  This proves \eqref{bls:finite-hard-lower}.

The proof covers arbitrary convex pieces, because it uses only a class's
vertices and their barycenter.  It does not use a prescribed variable
order, a branching rule, or an assumption that the planted vertex is
optimal.  The value comparison \eqref{app:bls-gap-paid} is what connects
the planted descent directions to the target at the actual integer
optimum.
